\begin{abstract}
Neural networks trained with empirical risk minimization on group-imbalanced data exhibit poor worst-group accuracy. Prior work established that minority groups occupy sharper regions of the loss landscape (sharpness ratio SR $> 1.0$), but the origin of this curvature asymmetry remained unclear. We investigate the \textbf{Differential Convergence Rate} hypothesis: majority groups converge faster during training, smoothing curvature in majority-relevant parameter directions while leaving minority directions sharp.

We present two key findings. First, we establish that \textbf{SR $\approx 1.0$ at random initialization} (SR$_0$ = 0.9999, 95\% CI [0.9953, 1.0046]), ruling out intrinsic data geometry as the source of curvature asymmetry. Second, we show that \textbf{gradient ratio decay temporally precedes sharpness ratio divergence} by $\tau = 3.67$ epochs (95\% CI [3.0, 4.0]). This positive lag indicates that majority convergence precedes minority curvature increase, consistent with a causal mechanism.

Our findings reveal a diagnostic window: gradient dynamics signal impending curvature divergence 3--4 epochs before it manifests. This temporal characterization opens the path from post-hoc worst-group corrections to early-training interventions targeting the root cause.
\end{abstract}
